In this section, we explore on the approaches and findings we uncovered while
tackling each of our theory objectives. More specifically,
while we have not achieved a complete proof for each objective,
we highlight several theorems
that may assist in their full resolution.
The core idea that we employed comes from following set of observations:
when working with \textit{PureCircuit}, one cannot fully work with binary logic, meaning
there is no trivial method for detecting $\bot$ values and removing them, as made clear by
Marino's continuity argument \cite{marino_GeneralTheoryMetastable_1981}.
We therefore make use of topological problems with the idea that resolutions of our input
corresponds to some set nearby points in a space. If all points abide some property
then we can use that to detect solutions.
This is an extension of the non-parsimonious reduction used by Deligkas
et al. \cite{deligkas_PureCircuitTightInapproximability_2024} when showing \textbf{PPAD}-hardness.


With respect to the first objective, we provide the following
list of theorems.
\begin{enumerate}
    \item $\textit{\#nD-StrongSperner} \subseteq \textit{\#HF-nD-StrongSperner}$ \ref{thm:hf-nd-ss-hard}
    \item $\textit{\#Hf-DiscreteStrongSperner} \subseteq \textit{\#PureCircuit}$ \ref{thm:hf-discrete-to-pure}
    \item $\textit{\#EndOfLine} \subseteq \textit{AntipodalStrongSperner}$ \ref{thm:parsimonious-2d-strong-eol-to-antipodalss}
\end{enumerate}
The problems mentioned are defined in the section \ref{sec:pc-plus-eol}, where
we also discuss their implications with respect to the current objective
and with \textbf{PPAD} as a whole

Following the same thought process for our second objective, we demonstated how to reduce \textit{SourceOrExcess}
to a topological problem. Specifically in section \ref{sec:sc-excess} we argue the following result
$$
    \textit{\#SourceOrExcess}(2,1) \subseteq \textit{\#3D-CountingStrongSperner} \text{ \ref{thm:counting-ss-excess}}
$$
Where we define a different combinatorial version of the \textit{3D-StrongSperner} and demonstrate
how it can be used for parsiomonious reductions between the $\textit{SourceOrExcess}(2,1)$ problem.


In the third section, we offer an combinatorial extension to the \textbf{PPAD} class
with regards to directed graphs and the number of known sources. Specifically
given Hollender et al.'s work \cite{hollender_ComplexityMultisourceVariants_2018}
on directed graphs with multiple known sources, we make
on observations as to the combinatorial nature of the problems that can be created
as well as important problems to analyse for a full \textit{PureCircuti} Cook-Levin type theorem.

Lastly we provide two supplementary findings with regards to the \textit{PureCircuit} reductions.
We first demonstrate how we can relate to \textit{Tarski} with \textit{Kleene theory} and
show how to create parsimonious reductions using \textit{PureCircuit}.
Second, we show how simplifications of the \textit{PureCircuit} problem can lead
to parsimonious reductions to the \textit{EndOfLine}. Since
these proofs to not offer any indication of the hardness of the problem,
we provide their full proofs in the appendix.